\documentclass{article}
\usepackage[a4paper, margin=0.7in]{geometry}
\usepackage[utf8]{inputenc}
\usepackage[T1]{fontenc}
\usepackage[french]{babel}
% Images
\usepackage{graphicx}
% Code packages
\usepackage{listings}
\usepackage{color}
 
\usepackage[most]{tcolorbox}
\definecolor{block-gray}{gray}{0.98}
\newtcolorbox{blockquote}{colback=block-gray,grow to right by=-1mm,grow to left by=-1mm,boxrule=0pt,boxsep=0pt,breakable}
 
% Defining the colors for the style
\definecolor{codegreen}{rgb}{0,0.6,0}
\definecolor{codegray}{rgb}{0.5,0.5,0.5}
\definecolor{codepurple}{rgb}{0.58,0,0.82}
\definecolor{backcolour}{rgb}{0.95,0.95,0.92}
 
% Defining the style
\lstdefinestyle{mystyle}{
    backgroundcolor=\color{backcolour},   
    commentstyle=\color{codegreen},
    keywordstyle=\color{magenta},
    numberstyle=\tiny\color{codegray},
    stringstyle=\color{codepurple},
    basicstyle=\footnotesize,
    breakatwhitespace=false,         
    breaklines=true,                 
    captionpos=b,                    
    keepspaces=true,                 
    numbers=left,                    
    numbersep=5pt,                  
    showspaces=false,                
    showstringspaces=false,
    showtabs=false,                  
    tabsize=2
}

% Setting the style
\lstset{style=mystyle}

% Listing -> Algorithme
\renewcommand{\lstlistingname}{Algorithme}
% List of Listings -> Algorithmes
\renewcommand{\lstlistlistingname}{\lstlistingname s}
% List of Figures -> Figures
\renewcommand{\listfigurename}{Figures}

\title{Correction du Partiel de PP, 2019-2020}
\author{Polio, Promo 2021}

\begin{document}

\maketitle

Durée: 1h30. Ce partiel est découpé en 5 exercices.

Correction par Polio (promo 2021), avec l'énoncé transcrit avant chaque réponse.

\section{Exercice I: Petit programme MPI}

\begin{blockquote}
Soit la section de code suivante exécutée par tous les processus d'un programme MPI (il y a au moins 2 processus):

\begin{lstlisting}[language=c]
int rang, som;

MPI_Comm_rank(MPI_COMM_WORLD, &rang);

som = 0;
if (rang == 0)
  som = 1;
som += 2;
if (rang == 1)
  som += 3;
\end{lstlisting}

Suite à l'exécution de ces lignes, quelle est la valeur de la variable \texttt{som} en fonction du rang MPI (justifiez brièvement) ?
\end{blockquote}

Tous les processus passent par \texttt{som = 0} et \texttt{som += 2}. Ainsi les processus dont le rang est supérieur ou égal à \texttt{2} auront pour valeur finale \texttt{som: 2}.

\textbf{Rang 0}: Passe par \texttt{som = 1} après le \texttt{som = 0}, la valeur finale sera donc \texttt{som: 3}.

\textbf{Rang 1}: Passe par \texttt{som += 3} après le \texttt{som += 2}, la valeur finale sera donc \texttt{som: 5}.

\section{Exercice II: Cherchez l'erreur}

\begin{blockquote}
Pour chacune des sections de code suivantes (numérotées \textbf{a} à \textbf{d}), dites:

\begin{itemize}
\item si elle fait l'object d'un risque de bloquage (``\emph{deadlock}'');
\item ou bien si les contenus des messages risquent d'être erronés;
\end{itemize}

Justifiez vos choix et apportez une correction pour que la section soit correcte.

Dans tout ce qui suit, on désigne par \texttt{rang} le rang du processus MPI et par \texttt{P} la taille du communicateur \texttt{MPI\_COMM\_WORLD}. On suppose que le programme a correctement initialisé MPI.
\end{blockquote}

\subsection{Section A}

\begin{blockquote}
On veut calculer le nombre total d'éléments nuls dans un tableau distribué.

\begin{lstlisting}[language=c]
int i, Nnuls = 0, Nnuls_tot;
/* tabl est un tableau distribue sur tous les processus et prealablement initialise
 * N = dimension locale du tableau
 */
for (i = 0; i < N; i++)
  if (tabl[i] == 0)
    Nnuls++;

if (Nnuls > 0) {
  MPI_Allreduce(&Nnuls, &Nnuls_tot, 1, MPI_INT, MPI_SUM MPI_COMM_WORLD);
  printf("Nb d'elements nuls dans le tableau distribue : %d\n", Nnuls_tot);
}
\end{lstlisting}
\end{blockquote}

\texttt{MPI\_Allreduce} est une opération collective: tous les processus compris dans \texttt{MPI\_COMM\_WORLD} doivent l'appeler.
Ici, seuls les processus ayant trouvé au moins un élément nul l'appellent, il y aura donc un \emph{deadlock}.

Pour corriger, il faut supprimer le \texttt{if (Nnuls $>$ 0)}.

\subsection{Section B}


\begin{blockquote}

Le processus 0 veut distribuer un entier différent à chacun des autres processus

\begin{lstlisting}[language=c]
MPI_Request req[P];
int entier;

if (rang == 0)
{
  for (i = 1; i < P; i++)
  {
    entier = i * i;
    MPI_Isend(&entier, 1, MPI_INT, i, 1000, MPI_COMM_WORLD, req + (i - 1));
  }
  MPI_Waitall(P - 1, req, MPI_STATUSES_IGNORE);
}
else
{
  MPI_Recv(&entier, 1, MPI_INT, 0, 1000, MPI_COMM_WORLD, MPI_STATUS_IGNORE);
}
\end{lstlisting}
\end{blockquote}

La variable \texttt{entier} est utilisée comme stockage pour chaque \texttt{MPI\_Isend(\&entier, [...])}. Si un de ces envois n'est pas complété avant que \texttt{entier} ne soit réécrit par le tour de boucle suivant, la valeur envoyée sera alors erronée.

Une correction consiste à déclarer un tableau \texttt{int entiers[P-1]} et à remplacer le contenu de la boucle par:

\begin{lstlisting}[language=c]
entiers[i - 1] = i * i;
MPI_Isend(entiers + (i - 1), 1, MPI_INT, i, 1000, MPI_COMM_WORLD, req + (i - 1));
\end{lstlisting}

Il faut toutefois laisser la variable \texttt{entier} afin qu'elle puisse être utilisée pour la réception pour les autres processus.

\subsection{Section C}


\begin{blockquote}

But: faire circuler un message (le contenu de \texttt{m}) dans un anneau initialisé par le processus \texttt{0}. Le message circule dans le sens croissants des rangs modulo \texttt{P}.

\begin{lstlisting}[language=c]
int droite, int gauche, m;

if (rang == 0)
{
  m = 0;
}

droite = (rang + 1) % P;
gauche = (rang - 1) % P;

MPI_Recv(&m, 1, MPI_INT, gauche, 1000, MPI_COMM_WORLD, MPI_STATUS_IGNORE);
MPI_Send(&m, 1, MPI_INT, droite, 1000, MPI_COMM_WORLD);
\end{lstlisting}
\end{blockquote}

\texttt{MPI\_Recv} (et \texttt{MPI\_Send}) est bloquante, hors tous les processus commencent par l'appeler: aucun \texttt{MPI\_Send} ne sera fait. Il y a donc un \emph{deadlock}.

Une correction possible consiste à spécialiser le code pour le processus 0:

\begin{lstlisting}[language=c]
if (rang == 0)
{
  MPI_Send(&m, 1, MPI_INT, droite, 1000, MPI_COMM_WORLD);
  MPI_Recv(&m, 1, MPI_INT, gauche, 1000, MPI_COMM_WORLD, MPI_STATUS_IGNORE);
}
else
{
  MPI_Recv(&m, 1, MPI_INT, gauche, 1000, MPI_COMM_WORLD, MPI_STATUS_IGNORE);
  MPI_Send(&m, 1, MPI_INT, droite, 1000, MPI_COMM_WORLD);
}
\end{lstlisting}

\subsection{Section D}


\begin{blockquote}

Le processus 0 veut envoyer le message \texttt{m = 100} au processus 1.

\begin{lstlisting}[language=c]
if (rang == 0)
{
  m = 100;
  MPI_Request req;
  MPI_Isend(&m, 1, MPI_INT, 1, 1000, MPI_COMM_WORLD, &req);
}
else if (rang == 1)
{
  MPI_Recv(&m, 1, MPI_INT, 0, 1000, MPI_COMM_WORLD, MPI_STATUS_IGNORE);
}
\end{lstlisting}
\end{blockquote}

\texttt{MPI\_Isend} est non bloquant: le message risque de ne jamais être envoyé et donc de bloquer le processus 1 indéfiniment.

Pour corriger il est possible d'ajouter la ligne \texttt{MPI\_Wait(\&req, MPI\_STATUS\_IGNORE);} en dessous de la ligne d'envoi.

\section{Exercice III: Quelle est la bonne valeur ?}

\begin{blockquote}

Soit la section de code suivante exécutée par tous les processus d'un programme MPI (il y a au moins 3 processus):

\begin{lstlisting}[language=c]
int rang, val, valA, valB;
MPI_Request req[2];

MPI_Comm_rank(MPI_COMM_WORLD, &rang);

if (rang == 0)
{
  MPI_Irecv(&valA, 1, MPI_INT, MPI_ANY_SOURCE, 1000, MPI_COMM_WORLD, req + 0);
  MPI_Irecv(&valB, 1, MPI_INT, MPI_ANY_SOURCE, 1111, MPI_COMM_WORLD, req + 1);
  MPI_Waitall(2, req, MPI_STATUSES_IGNORE);
}
else if (rang == 1)
{
  val = rang;
  MPI_Send(&val, 1, MPI_INT, 0, 1111, MPI_COMM_WORLD);
}
else if (rang == 2)
{
  val = rang;
  MPI_Send(&val, 1, MPI_INT, 0, 1000, MPI_COMM_WORLD);
}
\end{lstlisting}

\textbf{Question:} Suite à l'exécution de ces lignes, quelles sont les valeurs des variables \texttt{valA} et \texttt{valB} du processus de rang 0 (justifiez) ?
\end{blockquote}

La différenciation sur les réceptions du processus 0 se fait sur les tags (\texttt{1000} et \texttt{1111}).

Le processus 1 envoie sur le tag 1111 donc \texttt{valB = 1}.
Le processus 0 envoie sur le tag 1000 donc \texttt{valA = 2}.

\newpage\section{Exercice IV: Implémentation d'un \texttt{allgather}}


\begin{blockquote}

Voici le code de la fonction \texttt{allgather} et son programme principal: complétez-le.

\begin{lstlisting}[language=c]
void allgather(int in, int *out)
{
  int rang, P;

  /* 1. A COMPLETER: il faut affecter le nombre de processus dans `P` */
  /* 2. A COMPLETER: il faut affecter la variable `rang` */

  if (rang == 0)
  {
    MPI_Status sta;
    out[0] = in;
    for (int p = 1; p < P; p++)
    {
      int val;
      /* 3. A COMPLETER: reception dans `val` de la valeur envoyee par le processus de rang `p` */
      /* 4. A COMPLETER: mise a jour du tableau `out` */
    }
  }
  else
  {
    /* 5. A COMPLETER: envoi de `in` au processus 0 */
  }

  /* 6. A COMPLETER: le processus 0 diffuse `out` aux autres processus */
}

int main(int argc, char *argv[])
{
  int rang, P, *out;

  /* 7. A COMPLETER */
  /* 8. A COMPLETER: il faut affecter le nombre de processus dans `P` */
  /* 9. A COMPLETER: il faut affecter la variable `rang` */

  out = (int *)malloc(P * sizeof(int));
  allgather(rang, out);

  for (int p = 0; p < P; p++)
    printf("P%d, out[%d] = %d\n", rang, P, out[p]);
  free(out);

  /* 10. A COMPLETER */

  return 0;
}
\end{lstlisting}

Soit la fonction \texttt{void allgather(int in, int *out)} appelée en même temps par tous les processus MPI.
Elle collecte dans le tableau \texttt{out} le contenu de toutes les variables \texttt{in} de tous les processus \texttt{MPI}.
Tous les processus doivent récupérer ensuite les résultats. Le tableau \texttt{out} est préalable alloué au nombre total de processus.

La fonction \texttt{allgather} est implémentée dans le bloc de code ci-dessus, ainsi qu'un programme principal qui utilise cette fonction. Cependant, ce code n'est pas complet.

\textbf{Travail à faire}: compléter les lignes \texttt{\slash * À COMPLÉTER *\slash } pour que la fonction \texttt{allgather} et le programme principal soient corrects.
\end{blockquote}

\begin{itemize}
\item 1 et 8: \texttt{MPI\_Comm\_size(MPI\_COMM\_WORLD, \&P);}
\item 2 et 9: \texttt{MPI\_Comm\_rank(MPI\_COMM\_WORLD, \&rang);}
\item 3: \texttt{MPI\_Recv(\&val, 1, MPI\_INT, p, 1000, MPI\_COMM\_WORLD, \&sta);}
\item 4: \texttt{out[p] = val;}
\item 5: \texttt{MPI\_Send(\&in, 1, MPI\_INT, 0, 1000, MPI\_COMM\_WORLD);}
\item 6: \texttt{MPI\_Bcast(out, P, MPI\_INT, 0, MPI\_COMM\_WORLD);}
\item 7: \texttt{MPI\_Init(\&argc, \&argv);}
\item 8: \texttt{MPI\_Finalize();}
\end{itemize}

\newpage\section{Exercice V: Recouvrement}


\begin{blockquote}

Définition de la fonction \texttt{collect\_trait}:

\begin{lstlisting}[language=c]
void collect_trait(int *tab, int N) {
  int rank, P, k;
  MPI_Comm_rank(MPI_COMM_WORLD, &rank);
  MPI_Comm_size(MPI_COMM_WORLD, &P);

  if (rank == 0) {
    MPI_status sta;
    int *buf2 = (int *)malloc(N * sizeof(int));
    trait(tab, N);
    for (k = 1; k < P; k++) {
      MPI_Recv(buf2, N, MPI_INT, k, 1000, MPI_COMM_WORLD, &sta);
      trait(buf2, N);
    }
    free(buf2);
  }
  else {
    MPI_Send(tab, N, MPI_INT, 0, 1000, MPI_COMM_WORLD);
  }
}
\end{lstlisting}

On considère une fonction \texttt{collect\_trait} dont la définition se trouve ci-dessus. Elle est utilisée dans les conditions suivantes:

\begin{enumerate}
\item Cette fonction est appelée par tous les processus MPI (comme une opération collective);

\item Chaque processus a pré-alloué et rempli son propre tableau de \texttt{N} entiers;

\item La valeur \texttt{N} est identique pour tous les processus MPI;

\item La fonction \texttt{trait(int *tab, int N)} (appelée par le processus 0) effectue un traitement coûteux à partir des données du tableau \texttt{tab} passé en argument, sans modifier le contenu de ses données.

\end{enumerate}
\end{blockquote}

\begin{blockquote}
\textbf{Question 1:} Quelle est la modification à apporter à la réception pour que le processus 0 reçoive un message à partir de n'importe quel processus ?
\end{blockquote}

Il faut remplacer le \texttt{k} utilisé pour indiquer depuis quel processus réceptionner par \texttt{MPI\_ANY\_SOURCE}.

\begin{blockquote}

\textbf{Question 2:} Expliquez brièvement pourquoi la réception (par le processus 0) n'est pas recouverte par le traitement \texttt{trait}.
\end{blockquote}

\texttt{MPI\_Recv} est bloquant: la réception est donc forcément effectuée avant que le traitement des données reçues ne s'effectue.

\begin{blockquote}

\textbf{Question 3:} Modifiez la section de code exécutée par le processus 0 pour recouvrir la réception par le traitement \texttt{trait} tout en recevant les messages dans n'importe quel ordre.

\emph{Indication: utilisez deux buffers de \texttt{N} entiers, un pour recevoir les données, l'autre pour traiter les données déjà reçues.}
\end{blockquote}

\begin{lstlisting}[language=c]
MPI_status sta;
// Nouvelles variables utilisees pour le recouvrement
MPI_Request req;
int n;
int *buf2 = (int *)malloc(N * sizeof(int));
int *recvbuf = (int *)malloc(N * sizeof(int));
// Copie du tableau du rang 0 pour qu'il soit traite pendant la premiere reception
for (n = 0; n < P; n++) buf2[n] = tab[n];
// Recouvrement avec reception depuis toutes les sources
for (k = 1; k < P; k++) {
  MPI_Irecv(recvbuf, N, MPI_INT, MPI_ANY_SOURCE, 1000, MPI_COMM_WORLD, &req);
  // Traitement du tableau precedent pendant la reception du suivant
  trait(buf2, N);
  MPI_Wait(&req, &sta);
  for (n = 0; n < P; n++) buf2[n] = recvbuf[n]; // Copie
}
// Traitement du dernier tableau qui n'a pas ete fait apres la derniere copie: il faut donc le faire apres la boucle
trait(buf2, N);
free(buf2);
free(recvbuf); // Ne pas oublier de liberer le nouveau buffer
\end{lstlisting}

\end{document}
